% GENERATED FILE. Edit canonical lessons, then run npm run generate:books.
% canonical-source-hash: fnv1a64:fba3a4b1f5b9a755
% canonical-lessons: SA-C132-krsah, SA-C132-vakrah, SA-C132-rjuh, SA-C132-purnah, SA-C132-riktah

\chapter{Thin, Crooked, Straight, Full, Empty}
\label{ch:sa-thin-crooked-straight-full-empty}
\hlchaptermodality{132}

\begin{chapteropening}
\textbf{By the end of this chapter:} \emph{I can say thin, crooked, straight, full and empty.}

Complete the last lesson of chapter 132: I can say thin, crooked, straight, full and empty.
\end{chapteropening}

\section[\texorpdfstring{\sk{कृशः} (kṛśaḥ)}{कृशः (kṛśaḥ)}]{\sk{कृशः} (kṛśaḥ) — thin}

\label{lesson:SA-C132-krsah}



\begin{quote}
Before the new one: say the Sanskrit for clear, then the Sanskrit for thick.
\end{quote}

\subsection*{You'll want to know: \sk{कृशः}}
\textbf{\sk{कृशः}} — \emph{kṛśaḥ} — \textquotedblleft{}thin\textquotedblright{}.

\begin{grammarlens}[title={What you've built}]
One more word for this chapter.
\end{grammarlens}

\subsection*{Your turn}
\begin{itemize}
\raggedright
  \item \emph{Say it:} \emph{kṛśaḥ}
  \item \emph{Say it:} \emph{kṛśaḥ}, once more
  \item \emph{Say it:} say \emph{kṛśaḥ}
  \item \emph{Recall:} say the Sanskrit for clear, then the Sanskrit for thick, then say \emph{kṛśaḥ} again
\end{itemize}

\begin{tcolorbox}[breakable,colback=teal!4,colframe=teal!35!black,title={Before you move on}]
What does \sk{कृशः} mean? (\textquotedblleft{}Thin\textquotedblright{}.) Say it once more.
\end{tcolorbox}
\section[\texorpdfstring{\sk{वक्रः} (vakraḥ)}{वक्रः (vakraḥ)}]{\sk{वक्रः} (vakraḥ) — crooked}

\label{lesson:SA-C132-vakrah}



\begin{quote}
Before the new one: say the Sanskrit for thick, then the Sanskrit for thin.
\end{quote}

\subsection*{You'll want to know: \sk{वक्रः}}
\textbf{\sk{वक्रः}} — \emph{vakraḥ} — \textquotedblleft{}crooked\textquotedblright{}.

\begin{grammarlens}[title={What you've built}]
One more word for this chapter.
\end{grammarlens}

\subsection*{Your turn}
\begin{itemize}
\raggedright
  \item \emph{Say it:} \emph{vakraḥ}
  \item \emph{Say it:} \emph{vakraḥ}, once more
  \item \emph{Say it:} say \emph{vakraḥ}
  \item \emph{Recall:} say the Sanskrit for thick, then the Sanskrit for thin, then say \emph{vakraḥ} again
\end{itemize}

\begin{tcolorbox}[breakable,colback=teal!4,colframe=teal!35!black,title={Before you move on}]
What does \sk{वक्रः} mean? (\textquotedblleft{}Crooked\textquotedblright{}.) Say it once more.
\end{tcolorbox}
\section[\texorpdfstring{\sk{ऋजुः} (ṛjuḥ)}{ऋजुः (ṛjuḥ)}]{\sk{ऋजुः} (ṛjuḥ) — straight}

\label{lesson:SA-C132-rjuh}



\begin{quote}
Before the new one: say the Sanskrit for thin, then the Sanskrit for crooked.
\end{quote}

\subsection*{You'll want to know: \sk{ऋजुः}}
\textbf{\sk{ऋजुः}} — \emph{ṛjuḥ} — \textquotedblleft{}straight\textquotedblright{}.

\begin{grammarlens}[title={What you've built}]
One more word for this chapter.
\end{grammarlens}

\subsection*{Your turn}
\begin{itemize}
\raggedright
  \item \emph{Say it:} \emph{ṛjuḥ}
  \item \emph{Say it:} \emph{ṛjuḥ}, once more
  \item \emph{Say it:} say \emph{ṛjuḥ}
  \item \emph{Recall:} say the Sanskrit for thin, then the Sanskrit for crooked, then say \emph{ṛjuḥ} again
\end{itemize}

\begin{tcolorbox}[breakable,colback=teal!4,colframe=teal!35!black,title={Before you move on}]
What does \sk{ऋजुः} mean? (\textquotedblleft{}Straight\textquotedblright{}.) Say it once more.
\end{tcolorbox}
\section[\texorpdfstring{\sk{पूर्णः} (pūrṇaḥ)}{पूर्णः (pūrṇaḥ)}]{\sk{पूर्णः} (pūrṇaḥ) — full}

\label{lesson:SA-C132-purnah}



\begin{quote}
Before the new one: say the Sanskrit for crooked, then the Sanskrit for straight.
\end{quote}

\subsection*{You'll want to know: \sk{पूर्णः}}
\textbf{\sk{पूर्णः}} — \emph{pūrṇaḥ} — \textquotedblleft{}full\textquotedblright{}.

\begin{grammarlens}[title={What you've built}]
One more word for this chapter.
\end{grammarlens}

\subsection*{Your turn}
\begin{itemize}
\raggedright
  \item \emph{Say it:} \emph{pūrṇaḥ}
  \item \emph{Say it:} \emph{pūrṇaḥ}, once more
  \item \emph{Say it:} say \emph{pūrṇaḥ}
  \item \emph{Recall:} say the Sanskrit for crooked, then the Sanskrit for straight, then say \emph{pūrṇaḥ} again
\end{itemize}

\begin{tcolorbox}[breakable,colback=teal!4,colframe=teal!35!black,title={Before you move on}]
What does \sk{पूर्णः} mean? (\textquotedblleft{}Full\textquotedblright{}.) Say it once more.
\end{tcolorbox}
\section[\texorpdfstring{\sk{रिक्तः} (riktaḥ)}{रिक्तः (riktaḥ)}]{\sk{रिक्तः} (riktaḥ) — empty}

\label{lesson:SA-C132-riktah}



\begin{quote}
Before the new one: say the Sanskrit for straight, then the Sanskrit for full.
\end{quote}

\subsection*{You'll want to know: \sk{रिक्तः}}
\textbf{\sk{रिक्तः}} — \emph{riktaḥ} — \textquotedblleft{}empty\textquotedblright{}.

\begin{grammarlens}[title={What you've built}]
One more word for this chapter.
\end{grammarlens}

\subsection*{Your turn}
\begin{itemize}
\raggedright
  \item \emph{Say it:} \emph{riktaḥ}
  \item \emph{Say it:} \emph{riktaḥ}, once more
  \item \emph{Say it:} say \emph{riktaḥ}
  \item \emph{Recall:} say the Sanskrit for straight, then the Sanskrit for full, then say \emph{riktaḥ} again
\end{itemize}

\begin{tcolorbox}[breakable,colback=teal!4,colframe=teal!35!black,title={Before you move on}]
What does \sk{रिक्तः} mean? (\textquotedblleft{}Empty\textquotedblright{}.) Say it once more.
\end{tcolorbox}
